\documentclass[portrait, a1paper]{baposter}

\definecolor{light-blue}{rgb}{0, 0.85, 1}
\definecolor{dark-blue}{rgb}{0, 0.45 , 0.75}
\definecolor{white}{rgb}{1, 1, 1}

\usepackage[font=small,labelfont=bf]{caption}
\usepackage{float}
\usepackage{wrapfig}
\usepackage{graphicx}
\usepackage{caption}
\usepackage{subcaption}
\usepackage{setspace} 
\usepackage{relsize}
\newcommand{\imagedirectory}[1]{../images/{#1}}
\newcommand{\imagecommon}[1]{../../images/{#1}}

\begin{document}

%%% Setting Background Image %%%%%%%%%%%%%%%%%%%%%%%%%%%%%%%%%%%%%%%%%%%%%%%%%%
\background{
	\begin{tikzpicture}[remember picture,overlay]%
	\draw (current page.north west)+(-2em,2em) node[anchor=north west]
	{\includegraphics[height=1.1\textheight]{background}};
	\end{tikzpicture}
}

	\begin{poster}
	%poster option
	{
		%headershade=shade,
		headerColorOne=dark-blue,
		headerColorTwo=white,
		headerFontColor=white,
		headerfont=\bf,
		headerborder=open,
		headershape=rectangle,
		borderColor=light-blue,
		background=user,
		%bgColorOne=white,
		columns=3,
		eyecatcher=false,
		colspacing=6pt,
		linewidth=1pt,
		headerheight=0.085\textheight,
		textborder=rectangle,
		boxColorOne=white,
		 % Background color of the content boxes
	}
	{
    %eye catcher
	}
	%Title
	{
	\Large{Modelica Model for the youBot Manipulator} \vspace{0.1em}
	}
	%Author
	{
	\small{Rhama Dwiputra$^*$, Alexey Zakharov$^+$, Roustiam Chakirov$^*$, Erwin Prassler$^*$\\
	\scriptsize{
$^*${\{rhama.dwiputra, roustiam.chakirov, erwin.prassler\}@h-brs.de}, $^+${alexey.zhakarov@gmail.com}
		}
		}
	}
	%University Logo
	{
	\includegraphics[width=0.25\textwidth]{\imagecommon{logo-hbrs}}
	}

%---------------------------------------
% Background
%---------------------------------------

	\headerbox{Background}{name=background, column=0, row=0}{
		The \emph{youBot} is designed to serve as the reference platform for industry, research and education \cite{Bischoff2011}.
		Due to its frequent use for educational purpose and experiments, a model of the youBot is advantageous.  
	}
%---------------------------------------
% Approach
%---------------------------------------
	
	\headerbox{Approach}{name=Approach, column=1, row=0, span=2}{
\begin{wrapfigure}{r}{0.46\textwidth}
   \begin{center}
   \vspace{-5pt}
	\includegraphics[width=0.43\textwidth]{\imagedirectory{AbstractionLayer}}\\
 	\label{fig:abstraction-layers}
 	 \vspace{5pt}
 	 \smaller{Component-based development.}
 	 \end{center}
 \end{wrapfigure}	
	
The project emphasize on modularity and re-usability.
The manipulator model is developed with a library for its components. 
The model development is followed with a comparison test to observe the model accuracy.

}



%---------------------------------------
% Result and Evaluation
%---------------------------------------

	\headerbox{Result and Evaluation}{name=resultandevaluation, column=0, below=background, span=3}{
\begin{figure}[H]
   \vspace{-10pt}
	\begin{subfigure}[b]{0.72\textwidth}
	\centering
	\includegraphics[width=1\textwidth]{\imagedirectory{overview}}
	\label{fig:youBot-overview}	
	\end{subfigure}
	~
	\begin{subfigure}[b]{0.24\textwidth}
	\centering
	\includegraphics[width=1\textwidth]{\imagedirectory{parameters-and-evaluation}}
	\label{fig:youBot-overview}	
	\end{subfigure}\\
	\centering   
	\smaller{Model overview, parameter configuration and comparison test.}
\end{figure}
 \vspace{-10pt}
The benefit of having a model only holds true when the model is accurate \cite{Fritzon2004}.
The result of the comparison test shows that the model reflects the actual system within a reasonable deviation.

\begin{figure}[H]
   \vspace{-5pt}
	\begin{subfigure}[b]{0.3\textwidth}
	\centering
	\includegraphics[width=0.9\textwidth]{\imagedirectory{EvaluationErrorJointAngles}}
	\label{fig:youBot-overview}	
	\end{subfigure}
	~
	\begin{subfigure}[b]{0.3\textwidth}
	\centering
	\includegraphics[width=0.9\textwidth]{\imagedirectory{EvaluationJoint3}}
	\label{fig:EvaluationJoint3}	
	\end{subfigure}
	~
	\begin{subfigure}[b]{0.3\textwidth}
	\centering
	\includegraphics[width=0.9\textwidth]{\imagedirectory{EvaluationJoint5}}
	\label{fig:youBot-overview}	
	\end{subfigure}\\
	\centering  
	\smaller{Evaluation result}
\end{figure}

	
	}

%---------------------------------------
% Future work
%---------------------------------------

	\headerbox{Future Work}{name=FutureWork, column=0, below=resultandevaluation, span=3}{
Possible improvements is the development of a more accurate motor model and a more comprehensive evaluation of the manipulator model.
The manipulator model is planned to be tested with other Modelica tools (e.g., OpenModelica \cite{Fritzon2006}) and used for hardware-in-the-loop experiments.

	}

%---------------------------------------
% Acknowledgment
%---------------------------------------	
	
	\headerbox{Acknowledgement}{name=Acknowledgement, column=0, below=FutureWork}{
We gratefully acknowledge the continued support of the Bonn-Rhein-Sieg University of Applied Sciences and its academic staffs for the help and support in this project. 
	}

%---------------------------------------
% Reference
%---------------------------------------	
	
	\headerbox{Reference}{name=Reference, column=1, below=FutureWork, span=2}{
	\smaller{
\vspace{-0.4em} 										% Save some space at the beginning
\bibliographystyle{plain}							% Use plain style
\renewcommand{\section}[2]{\vskip 0.05em}		% Omit "References" title

\begin{thebibliography}{1}							% Simple bibliography with widest label of 1
\itemsep=-0.01em										% Save space between the separation
\setlength{\baselineskip}{0.4em}					% Save space with longer lines
\bibitem{Bischoff2011} R. Bischoff, U. Huggenberger, and E. Prassler, \emph{Kuka youbot - a mobile manipulator for research and education}, in \emph{IEEE Int. Conf. on Robotics and Automat. (ICRA)}, pp. 1-4, May 2011.
\bibitem{Fritzon2004} P. Fritzson, \emph{Principles of Object-Oriented Modeling and Simulation With Modelica 2.1.} IEEE Press, 2004.
\bibitem{Fritzon2006} P. Fritzson et al., \emph{OpenModelica - A Free Open-Source Environment for System Modeling, Simulation, and Teaching} in \emph{IEEE Int. Conf. on Control Appl.}, 2006.
\end{thebibliography}

	}
	}
		
		
		
	\end{poster}
\end{document}